% UBT-AI-PROVENANCE-BEGIN
% schema: ubt-ai-provenance/v1
% tier: C_working
% ai_assistance: disclosed
% human_review: risk-based
% editorial_responsibility: Ing. David Jaroš
% policy: ../../../AI_PROVENANCE.md
% notice: Working material; exhaustive human review is not claimed.
% UBT-AI-PROVENANCE-END

\chapter{Theta, Mellin, zeta, paths and prime structure}
\label{ch:theta-mellin-feynman}

\section{Why this chapter exists}

This chapter is deliberately broader than a canonical UBT claim document.  Its purpose is
pedagogical: to place several classical mathematical structures next to the reduced UBT
theta bridge so that their relations can be seen in one place.  A result does\emph{not}
become a UBT theorem merely because it is useful inside UBT.

We therefore use three status labels in the prose:
\begin{itemize}
  \item\textbf{Classical}: standard mathematics or standard quantum mechanics;
  \item\textbf{Bridge}: an exact statement about the reduced theta model used to connect
        standard mathematics to UBT notation;
  \item\textbf{UBT open}: a statement that would have to be derived from the canonical
        UBT field equations or action before it could become a physical UBT claim.
\end{itemize}

The central reduced object is
\begin{equation}
  S(t,\psi)
  = \sum_{n\ge 1} a_n\,e^{-\pi\psi n^2}e^{i\pi t n^2}
  = \sum_{n\ge1}a_n e^{-\pi(\psi-it)n^2}.
  \label{eq:reduced-theta-amplitude}
\end{equation}
It is the same reduced complex-time kernel used in
\path{canonical/bridges/theta_complex_time_bridge.tex}.  In this chapter it is a
\textbf{bridge model}; the canonical UBT master field $\Theta(q,\tau)$ is not assumed to
be a Jacobi theta function.

\paragraph{Sign convention used in related RH notes.}
Some earlier UBT/RH scratch work used a positive heat variable $y=-\psi$ and therefore
wrote the theta parameter as $t-i\psi=t+iy$ with convergence on the corresponding
$y>0$ half-line.  The present chapter writes the damping variable directly as $\psi>0$
in Eq.~\eqref{eq:reduced-theta-amplitude}.  The transform identities are identical after
the relabelling $y=-\psi$; the sign convention must be declared before a physical UBT
interpretation is attached to either half-line.

\section{One spectrum, several representations}

\subsection{Orthogonality of modes versus nonorthogonality of theta templates}

For fixed $\psi$, the elementary time modes in Eq.~\eqref{eq:reduced-theta-amplitude} are
\begin{equation}
  e^{i\pi n^2 t}.
\end{equation}
On the interval $t\in[0,2]$ they are orthogonal:
\begin{equation}
 \int_0^2 e^{i\pi(n^2-m^2)t}\,dt
 = 2\,\delta_{nm},
 \qquad n,m\ge1.
 \label{eq:quadratic-mode-orthogonality}
\end{equation}
The important distinction is that the\emph{complete templates}
$S_\psi(t):=S(t,\psi)$ at different values of $\psi$ are generally not orthogonal.

\begin{theorem}[Classical Gram kernel identity]
For the inner product
\begin{equation}
  \langle f,g\rangle:=\frac12\int_0^2 f(t)\overline{g(t)}\,dt,
\end{equation}
and amplitudes $S_\psi$ of Eq.~\eqref{eq:reduced-theta-amplitude},
\begin{equation}
  G(\psi,\phi)
  :=\langle S_\psi,S_\phi\rangle
  =\sum_{n\ge1}|a_n|^2e^{-\pi(\psi+\phi)n^2}.
  \label{eq:gram-general}
\end{equation}
For $a_n\equiv1$,
\begin{equation}
  G(\psi,\phi)
  =\frac{\vartheta_3(0\mid i(\psi+\phi))-1}{2}.
  \label{eq:gram-theta}
\end{equation}
\end{theorem}

\begin{proof}
Insert the two sums in the inner product.  Equation~\eqref{eq:quadratic-mode-orthogonality}
eliminates all $m\neq n$ cross terms, leaving
\[
  \sum_n |a_n|^2 e^{-\pi\psi n^2}e^{-\pi\phi n^2}.
\]
For unit weights, the standard identity
$\vartheta_3(0\mid iu)=1+2\sum_{n\ge1}e^{-\pi u n^2}$ gives
Eq.~\eqref{eq:gram-theta}.
\end{proof}

\paragraph{Practical consequence.}
Template non-orthogonality is not an obstacle to detection.  If an observed signal is
$f(t)$, a normalized matched-filter score is
\begin{equation}
  \rho(\psi,\Delta)
  =\frac{\left|\int f(t)\overline{S(t-\Delta,\psi)}\,dt\right|^2}
  {\|f\|^2\,\|S_\psi\|^2}.
  \label{eq:theta-matched-filter}
\end{equation}
The best single template is obtained by maximizing $\rho$ over $(\psi,\Delta)$.
Orthogonality becomes important when many templates are fitted simultaneously; then the
Gram matrix must be inverted or regularized, for example through $G+\lambda I$.

\subsection{Real-time Fourier transform}

Take the Fourier convention
\begin{equation}
  \widehat f(\omega)=\int_{-\infty}^{\infty}f(t)e^{-i\omega t}\,dt.
\end{equation}
Formally, or distributionally,
\begin{equation}
  \widehat S(\omega,\psi)
  =2\pi\sum_{n\ge1}a_n e^{-\pi\psi n^2}
       \delta(\omega-\pi n^2).
  \label{eq:fourier-delta-spectrum}
\end{equation}
Thus the quadratic mode spectrum is already a Dirac comb in real-time frequency; no
mode-by-mode demodulation is needed to create the delta lines.

\subsection{Laplace transform in imaginary time}

For $\Re z$ sufficiently large,
\begin{align}
  \mathcal L_\psi[S](z;t)
   &=\int_0^\infty e^{-z\psi}S(t,\psi)\,d\psi \\
   &=\sum_{n\ge1}
      \frac{a_n e^{i\pi t n^2}}{z+\pi n^2}.
  \label{eq:laplace-resolvent}
\end{align}
The same eigenvalues that gave Fourier delta lines now appear as resolvent poles
$z=-\pi n^2$.

\subsection{Mellin's transformation in imaginary time}

For $\Re s$ in the absolute-convergence half-plane,
\begin{align}
  \mathcal M_\psi[S](s;t)
  &:=\int_0^\infty \psi^{s-1}S(t,\psi)\,d\psi \\
  &=\Gamma(s)\pi^{-s}
    \sum_{n\ge1}\frac{a_n e^{i\pi t n^2}}{n^{2s}}.
  \label{eq:mellin-time-twisted}
\end{align}
The identity follows from
\begin{equation}
 \int_0^\infty \psi^{s-1}e^{-\pi n^2\psi}\,d\psi
 =\Gamma(s)(\pi n^2)^{-s}.
\end{equation}

For the Gram kernel with $a_n\equiv1$ and $u=\psi+\phi$,
\begin{equation}
 \int_0^\infty u^{s-1}G(u)\,du
 =\Gamma(s)\pi^{-s}\zeta(2s).
 \label{eq:gram-mellin-zeta}
\end{equation}
This is classical theta--Mellin mathematics.  Its relevance here is organizational:
the same reduced theta bank that causes non-orthogonal templates is also a heat trace
whose Mellin representation is a spectral zeta function.

\paragraph{The spectral triangle.}
Equations~\eqref{eq:fourier-delta-spectrum},\eqref{eq:laplace-resolvent}, and
\eqref{eq:mellin-time-twisted} give three views of the same quadratic spectrum:
\begin{equation}
 \boxed{
 \begin{array}{ccl}
   \mathcal F_t   &\longrightarrow& \delta(\omega-\pi n^2),\\[2mm]
   \mathcal L_\psi&\longrightarrow& (z+\pi n^2)^{-1},\\[2mm]
   \mathcal M_\psi&\longrightarrow& \Gamma(s)\pi^{-s}n^{-2s}.
 \end{array}}
 \label{eq:spectral-triangle}
\end{equation}

\section{Mellin is the Fourier transform in logarithmic coordinate}

Let
\begin{equation}
  M_f(s)=\int_0^\infty f(x)x^{s-1}\,dx,
  \qquad s=\sigma+i\omega.
\end{equation}
Set $x=e^u$, so $dx=e^u du$.  Then
\begin{equation}
  M_f(\sigma+i\omega)
  =\int_{-\infty}^{\infty}
   f(e^u)e^{\sigma u}e^{i\omega u}\,du.
  \label{eq:mellin-log-fourier}
\end{equation}
Apart from the conventional sign of the Fourier phase, this is a Fourier transform in
$u=\ln x$ of the weighted signal $f(e^u)e^{\sigma u}$.  Thus Mellin analysis is the
natural harmonic analysis for multiplicative scaling, whereas ordinary Fourier analysis
is natural for additive translation.

This observation is especially useful when interpreting zeta functions: zeta's vertical
variable $\Im s$ is a log-scale frequency.

\section{Current cut and time-phased family of zeta functions}
\label{sec:time-twisted-zeta}

\paragraph{UBT coordinate convention.}
In the working complex-time interpretation used in this research thread, the present
slice is chosen as $t=0$.  This is a UBT coordinate convention, not a statement that the
Riemann zeta function intrinsically knows a physical ``present.''

For unit weights $a_n\equiv1$, define the quadratic time twist
\begin{equation}
  Z_\Theta(s;t)
  :=\sum_{n=1}^{\infty}
    \frac{e^{i\pi t n^2}}{n^{2s}}.
  \label{eq:time-twisted-zeta}
\end{equation}
Then Eq.~\eqref{eq:mellin-time-twisted} becomes
\begin{equation}
 \boxed{
  \mathcal M_\psi[S](s;t)
  =\Gamma(s)\pi^{-s}Z_\Theta(s;t).}
 \label{eq:mellin-Ztheta}
\end{equation}
At the present slice,
\begin{equation}
 \boxed{Z_\Theta(s;0)=\zeta(2s).}
 \label{eq:present-zeta}
\end{equation}
Thus ordinary zeta is one member---the $t=0$ member---of a natural family generated by
the same reduced complex-time theta kernel.

\begin{remark}
The exponential factor in Eq.~\eqref{eq:time-twisted-zeta} is an\emph{additive quadratic
twist}.  For generic $t$ the resulting Dirichlet series does not have the ordinary Euler
product.  One must not infer prime-factor dynamics merely from the existence of the
$t=0$ Euler product.
\end{remark}

\subsection{Simple special times}

The family is periodic under $t\mapsto t+2$ because $n^2\in\mathbb Z$:
\begin{equation}
 Z_\Theta(s;t+2)=Z_\Theta(s;t).
\end{equation}
At $t=1$,
\begin{align}
 Z_\Theta(s;1)
 &=\sum_{n\ge1}\frac{(-1)^{n^2}}{n^{2s}}
 =\sum_{n\ge1}\frac{(-1)^n}{n^{2s}}\\
 &=-\eta(2s)
 =-\bigl(1-2^{1-2s}\bigr)\zeta(2s),
 \label{eq:t1-eta}
\end{align}
where $\eta$ is the Dirichlet eta function.

\subsection{Rational times: decomposition using the Hurwitz zeta function}

Let $t=a/q$ with integers $a$ and $q>0$.  The coefficient
\begin{equation}
 c_n=e^{i\pi a n^2/q}
\end{equation}
is periodic with period $M=2q$(sometimes the true period is smaller).  Splitting the
Dirichlet series into residue classes gives the exact classical identity
\begin{equation}
 \boxed{
 Z_\Theta\!\left(s;\frac aq\right)
 =M^{-2s}\sum_{r=1}^{M}
 e^{i\pi a r^2/q}
 \zeta\!\left(2s,\frac rM\right),
 \qquad M=2q,}
 \label{eq:rational-hurwitz}
\end{equation}
where $\zeta(s,\alpha)$ is the Hurwitz zeta function.

\begin{proof}
Write every $n\ge1$ uniquely as $n=mM+r$ with $m\ge0$ and $1\le r\le M$.
Periodicity gives $c_{mM+r}=c_r$, and
\begin{align*}
 Z_\Theta(s;a/q)
 &=\sum_{r=1}^{M}c_r\sum_{m=0}^{\infty}(mM+r)^{-2s}\\
 &=M^{-2s}\sum_{r=1}^{M}c_r
   \zeta\!\left(2s,\frac rM\right).
\end{align*}
\end{proof}

This makes the ``alternative zeta functions'' generated away from $t=0$ precise.  The
ordinary zeta, eta, Hurwitz zeta combinations, and quadratic Gauss sums are not unrelated
objects here: they are different slices or representations of one quadratic-phase family.

\section{Rational Revivals and Quadratic Gaussian Sums}
\label{sec:gauss-revivals}

At rational $t=a/q$, the phase $e^{i\pi a n^2/q}$ depends only on a finite residue class.
The corresponding finite quadratic sums are Gauss sums.  A common normalization is
\begin{equation}
 g(a,q)=\sum_{r=0}^{q-1}e^{2\pi i a r^2/q}.
 \label{eq:gauss-sum}
\end{equation}
For coprime moduli $q_1,q_2$, the Chinese remainder theorem gives, with the
normalization in Eq.~\eqref{eq:gauss-sum},
\begin{equation}
 g(a,q_1q_2)=g(aq_2,q_1)\,g(aq_1,q_2),
 \qquad \gcd(q_1,q_2)=1.
 \label{eq:gauss-crt}
\end{equation}
Thus composite denominators factor into coprime pieces, and prime powers are the
irreducible arithmetic building blocks of rational-time revival structure.

This is a more promising place to search for a\emph{dynamical} origin of UBT prime
sectors than simply observing the Euler product of $\zeta(2s)$, because the Gauss-sum
arithmetic is carried directly by the quadratic theta phase.

\section{Sum of theta modes and Feynman sum of windings}
\label{sec:theta-feynman}

The theta/path relation is exact in a controlled textbook example: a free particle of
mass $m$ moving on a circle of circumference $L$.

The momentum eigenfunctions and energies are
\begin{equation}
 \phi_n(x)=L^{-1/2}e^{2\pi i n x/L},
 \qquad
 E_n=\frac{\hbar^2}{2m}\left(\frac{2\pi n}{L}\right)^2.
\end{equation}
Hence the real-time propagator is
\begin{equation}
 K(x,x';T)
 =\frac1L\sum_{n\in\mathbb Z}
 \exp\!\left(\frac{2\pi i n(x-x')}{L}
             -\frac{iE_nT}{\hbar}\right).
 \label{eq:circle-mode-propagator}
\end{equation}
This is a Jacobi-theta mode sum.  Poisson summation transforms it into
\begin{equation}
 \boxed{
 K(x,x';T)
 =\sqrt{\frac{m}{2\pi i\hbar T}}
  \sum_{w\in\mathbb Z}
  \exp\!\left[
   \frac{i m(x-x'+wL)^2}{2\hbar T}
  \right].}
 \label{eq:circle-winding-propagator}
\end{equation}
The integer $w$ labels winding sectors.  The exponent is the classical free-particle
action for the lifted path from $x'$ to $x+wL$:
\begin{equation}
 S_w^{\rm cl}=\frac{m(x-x'+wL)^2}{2T}.
\end{equation}
Therefore the same propagator admits two exact representations:
\begin{equation}
 \boxed{
 K_\Sigma
 \quad\Longleftrightarrow\quad
 K_{\mathbb Z}}
 \label{eq:mode-winding-duality}
\end{equation}
Here $K_\Sigma$ denotes the spectral mode sum and $K_{\mathbb Z}$ the sum over
winding sectors,
with Poisson/Jacobi transformation as the bridge.

\paragraph{Important qualification.}
The classical trajectory is selected by\emph{stationary} action,
$\delta S=0$, not necessarily minimum action.  Several stationary paths are ordinary
quantum mechanics; by themselves they do not imply a multiverse interpretation.

\section{Where prime numbers come in and where they don't}

\subsection{Euler's product on an unphased section}

At $t=0$ and unit weights,
\begin{equation}
 Z_\Theta(s;0)=\zeta(2s)
 =\prod_{p}\left(1-p^{-2s}\right)^{-1},
 \qquad \Re s>\frac12.
 \label{eq:zeta-euler-product}
\end{equation}
This is a profound change of representation from an additive sum over all positive
integers to a multiplicative product over primes.  It is nevertheless a classical
property of the integers, not by itself evidence that UBT dynamics selects prime modes.
For general weights $a_n$, the Dirichlet series
$\sum |a_n|^2n^{-2s}$ has an Euler product only when the coefficients satisfy the
appropriate multiplicativity conditions.

\subsection{Grids of poles of local Euler factors}

A single formal local factor
\begin{equation}
 L_p(s)=\left(1-p^{-2s}\right)^{-1}
\end{equation}
has poles when $p^{-2s}=1$.  Writing $s=\sigma+i\omega$ gives
\begin{equation}
 \sigma=0,
 \qquad
 \omega_{p,k}=\frac{\pi k}{\ln p},
 \qquad k\in\mathbb Z.
 \label{eq:local-prime-poles}
\end{equation}
For distinct primes $p\neq q$, nonzero pole frequencies cannot coincide.  If
$k/\ln p=\ell/\ln q$ with nonzero integers $k,\ell$, then
$q^k=p^\ell$, contradicting unique factorization.  All local factors share the trivial
$k=0$ location.

These are poles of the\emph{individual formal Euler factors}.  They lie outside the
Euler-product convergence half-plane and must not be confused with poles of the
analytically continued global $\zeta(2s)$, whose pole is at $s=1/2$.

If $N(T;P)$ counts positive local-factor pole frequencies up to height $T$ for primes
$p\le P$, then
\begin{align}
 N(T;P)
 &=\sum_{p\le P}\left\lfloor\frac{T\ln p}{\pi}\right\rfloor\\
 &=\frac{T}{\pi}\vartheta(P)+O(\pi(P)),
 \label{eq:local-pole-count}
\end{align}
where
\begin{equation}
 \vartheta(P)=\sum_{p\le P}\ln p
\end{equation}
is Chebyshev's theta function.  The prime number theorem implies
$\vartheta(P)\sim P$.

\subsection{Vertical Fourier spectrum and powers of prime numbers}

At fixed $\sigma>1/2$,
\begin{equation}
 \zeta\bigl(2(\sigma+i\omega)\bigr)
 =\sum_{n\ge1}n^{-2\sigma}e^{-i(2\ln n)\omega}.
 \label{eq:zeta-vertical-fourier}
\end{equation}
Thus Fourier transformation in $\omega$ produces spectral lines at
\begin{equation}
 \xi_n=2\ln n.
\end{equation}
To isolate primes and prime powers, use the logarithmic derivative:
\begin{equation}
 -\frac{\zeta'(z)}{\zeta(z)}
 =\sum_{n\ge1}\frac{\Lambda(n)}{n^z}
 =\sum_p\sum_{k\ge1}(\ln p)p^{-kz},
 \qquad \Re z>1.
 \label{eq:log-derivative-von-mangoldt}
\end{equation}
With $z=2(\sigma+i\omega)$,
\begin{equation}
 -\frac{\zeta'\!\left(2(\sigma+i\omega)\right)}
        {\zeta\!\left(2(\sigma+i\omega)\right)}
 =\sum_p\sum_{k\ge1}
   (\ln p)p^{-2k\sigma}
   e^{-i(2k\ln p)\omega}.
 \label{eq:prime-power-fourier-comb}
\end{equation}
The prime-power frequency comb therefore sits at
\begin{equation}
 \boxed{\xi_{p,k}=2k\ln p.}
 \label{eq:prime-power-frequencies}
\end{equation}
Distinct prime-power lines coincide only when the prime powers themselves are identical.

\subsection{Counting the primes and zero frequencies of the zeta function}

The weighted prime-power counting function
\begin{equation}
 \psi_C(x)=\sum_{n\le x}\Lambda(n)
\end{equation}
is related by the classical explicit formula to the nontrivial zeros
$\rho=\beta+i\gamma$ of $\zeta$.  Schematically,
\begin{equation}
 \psi_C(x)=x-\sum_\rho\frac{x^\rho}{\rho}+\cdots.
 \label{eq:explicit-formula-schematic}
\end{equation}
Writing $x=e^u$ turns each zero contribution into
\begin{equation}
 x^\rho=e^{\beta u}e^{i\gamma u}.
\end{equation}
Thus the imaginary parts $\gamma$ of zeta zeros act as oscillation frequencies in the
logarithmic variable $u=\ln x$.  This is the classical primes--zeros spectral bridge.

\section{What this means for UBT prime sectors}

Two distinct mechanisms must not be conflated.

\paragraph{Route A: trace/orbit route.}
Trace-formula analogies suggest that primes should correspond to primitive periodic
objects whose repeated traversals generate prime powers.  For the ordinary free particle
on $S^1$, however, the primitive-orbit length spectrum is not $\{\ln p\}$.  Therefore
\begin{quote}
\textbf{GAP-THETA-PRIME-1:} derive from a genuine UBT operator a primitive-orbit or
spectral structure whose natural primitive data are $\ln p$, or do not claim that the
Euler product dynamically derives UBT prime sectors.
\end{quote}

\paragraph{Route B: rational-revival route.}
The quadratic theta dynamics itself has arithmetic structure at rational times.  Gauss
sums factor through coprime denominators, making prime powers natural building blocks.
This motivates
\begin{quote}
\textbf{GAP-THETA-PRIME-2:} rewrite the existing UBT $\alpha$ prime-sector construction
in revival/winding variables and test whether its distinguished primes arise from
rational-time Gauss-sum structure rather than from a post-selected list.
\end{quote}

These routes are compatible as research directions but are not yet known to be the same
mechanism.

\section{Operator question behind the whole picture}

Equation~\eqref{eq:reduced-theta-amplitude} can be written formally by introducing
\begin{equation}
 H_0|n\rangle=\pi n^2|n\rangle.
\end{equation}
Then its mode evolution is generated by
\begin{equation}
 U_0(t,\psi)=e^{-\psi H_0}e^{+itH_0}.
 \label{eq:formal-complex-propagator}
\end{equation}
This makes the reduced theta amplitude genuinely propagator-like.  The decisive UBT
question is stronger:
\begin{quote}
\textbf{GAP-THETA-PROP:} derive an operator $H_\Theta$ from the canonical UBT action or
field equations such that an appropriate kernel or matrix element of
$e^{-\psi H_\Theta}e^{\pm itH_\Theta}$ yields the theta structure used here.
\end{quote}
Until this gap is closed, the mode/winding, Mellin/zeta, and revival/prime relations are
rigorous mathematics of the reduced bridge model, not a derivation of UBT dynamics.

\section{Big picture map}

The useful synthesis is therefore
\begin{equation}
\boxed{
\begin{array}{ccccc}
 & \mathcal K_\Theta & & &\\
 & \displaystyle\sum_n a_n e^{-\pi\psi n^2}e^{i\pi t n^2} & & &\\[2mm]
 \swarrow & \downarrow & \searrow & &\\[-1mm]
 \mathcal F_t/\mathcal L_\psi & \mathcal M_\psi & \mathcal P & &\\
 \mathcal D & Z_\Theta(s;t) & \mathcal W & &\\
 & \downarrow\ t=0 & \downarrow\ t\in\mathbb Q & &\\
 & \zeta(2s) & \mathcal G_{\mathbb Q} & &\\
 & \downarrow & & &\\
 & \mathcal E_p & & &
\end{array}}
\end{equation}
In the diagram, $\mathcal K_\Theta$ is the reduced theta kernel,
$\mathcal P$ is the Poisson/Jacobi transform, $\mathcal D$ denotes the
delta-line/resolvent branch, $\mathcal W$ the winding/path sum,
$\mathcal G_{\mathbb Q}$ rational-time Gauss arithmetic, and $\mathcal E_p$
the Euler-product, prime-power, and zero branch.
The classical arrows are exact.  The research task is to determine which of them survive
when the reduced kernel is replaced by a canonical UBT operator and whether either prime
route explains the existing UBT $\alpha$ sectors.
